\documentclass[aspectratio=169]{beamer}
\usetheme{metropolis}

\usepackage[utf8]{inputenc}
\usepackage[spanish]{babel}
\usepackage{graphicx}
\usepackage{xcolor}
\usepackage{tikz}
\usetikzlibrary{positioning, arrows.meta, babel}
\usepackage{tcolorbox}
\tcbuselibrary{skins, breakable}
\usepackage{fontawesome5}
\usepackage{amsmath}

% ─────────────────────────────────────────────────────────────
%  PALETA
% ─────────────────────────────────────────────────────────────
\definecolor{mainColor}{HTML}{00a499}
\definecolor{accentColor}{HTML}{EA7600}
\definecolor{bgLight}{HTML}{E6F7F6}
\definecolor{codeGray}{HTML}{2B2B2B}

% ─────────────────────────────────────────────────────────────
%  METROPOLIS
% ─────────────────────────────────────────────────────────────
\setbeamercolor{normal text}{fg=codeGray, bg=white}
\setbeamercolor{frametitle}{bg=mainColor, fg=white}
\setbeamercolor{alerted text}{fg=accentColor}
\setbeamercolor{progress bar}{fg=accentColor, bg=mainColor!25}
\setbeamercolor{block title}{bg=mainColor, fg=white}
\setbeamercolor{block body}{bg=bgLight, fg=codeGray}
\setbeamercolor{block title alerted}{bg=accentColor, fg=white}
\setbeamercolor{block body alerted}{bg=accentColor!10, fg=codeGray}
\metroset{progressbar=frametitle, sectionpage=none, numbering=fraction}
\usepackage{hyperref}
% ─────────────────────────────────────────────────────────────
%  CAJA DE CONCEPTO
% ─────────────────────────────────────────────────────────────
\newtcolorbox{conceptbox}[1]{
  colback=bgLight, colframe=mainColor,
  title={\small #1}, fonttitle=\bfseries,
  coltitle=white, colbacktitle=mainColor,
  arc=4pt, boxrule=0.6pt,
  left=4pt, right=4pt, top=2pt, bottom=2pt
}

\newcommand{\tech}[1]{\texttt{\textcolor{mainColor}{#1}}}

\hypersetup{colorlinks=true, urlcolor=mainColor}

% ─────────────────────────────────────────────────────────────
%  PORTADA
% ─────────────────────────────────────────────────────────────
\title{Modelos de Inteligencia Artificial en Negocios\\{\normalsize |}\\ Inteligencia Artificial Aplicada a la Industria\\{\tiny Clase 1: Introducción y Primeros Pasos}}
\subtitle{{\small Semestre 1, 2026 · Universidad de Santiago de Chile}}
\author{\href{https://www.linkedin.com/in/byron-caices/}{Byron Caices Lima} \\ \small Departamento de Ingeniería Industrial}
\date{\today}

\begin{document}

\maketitle

% =============================================================
%  SECCIÓN 1: BIENVENIDA Y CONTEXTO DEL CURSO
% =============================================================

% ─────────────────────────────────────────────────────────────
%  1. BIENVENIDA
% ─────────────────────────────────────────────────────────────
\begin{frame}{Bienvenidos al Curso}
  \begin{columns}[T]
    \begin{column}{0.55\textwidth}
      \small
      \textbf{¿De qué se trata este curso?}
      \begin{itemize}\setlength\itemsep{4pt}
        \item Capacitar a futuros ingenieros civiles industriales en los
              \textbf{fundamentos estratégicos} de la IA.
        \item No es un curso de programación: el foco está en
              \textbf{comprender, evaluar y aplicar} soluciones de IA
              en contextos industriales.
        \item Integraremos dimensiones \textbf{éticas, tecnológicas
              y organizacionales}.
      \end{itemize}
    \end{column}
    \begin{column}{0.42\textwidth}
      \begin{conceptbox}{\faRocket\enspace Objetivo General}
        \small
        Proporcionar herramientas conceptuales para comprender y aplicar
        la IA como \textbf{herramienta de gestión estratégica} y toma
        de decisiones en la industria.
      \end{conceptbox}
    \end{column}
  \end{columns}
\end{frame}

% ─────────────────────────────────────────────────────────────
%  2. ESTRUCTURA DEL CURSO
% ─────────────────────────────────────────────────────────────
\begin{frame}{Estructura del Curso: 16 Semanas}
  \small
  \begin{columns}[T]
    \begin{column}{0.48\textwidth}
      \begin{block}{\faBookOpen\enspace Unidades 1--5}
        \begin{enumerate}\setlength\itemsep{2pt}
          \item Introducción a la IA
          \item IA Estratégica
          \item Datos para IA
          \item Aprendizaje Automático
          \item Modelos Neuronales
        \end{enumerate}
      \end{block}
    \end{column}
    \begin{column}{0.48\textwidth}
      \begin{block}{\faCogs\enspace Unidades 6--10}
        \begin{enumerate}\setcounter{enumi}{5}\setlength\itemsep{2pt}
          \item Arquitecturas Avanzadas
          \item Evaluación de Soluciones
          \item Ética y Regulación
          \item Oportunidades Emergentes
          \item Proyecto Final
        \end{enumerate}
      \end{block}
    \end{column}
  \end{columns}
  \vspace{0.3cm}
  \begin{alertblock}{\faBalanceScale\enspace Evaluación}
    \centering\small
    \textbf{Informes de Laboratorio} (50\,\%) \quad+\quad
    \textbf{Proyecto Final Grupal} (50\,\%)
  \end{alertblock}
\end{frame}

% ─────────────────────────────────────────────────────────────
%  3. METODOLOGÍA
% ─────────────────────────────────────────────────────────────
\begin{frame}{Metodología: ¿Cómo Trabajaremos?}
  \begin{columns}[T]
    \begin{column}{0.48\textwidth}
      \begin{conceptbox}{\faChalkboardTeacher\enspace Clases Teóricas}
        \small
        Sesiones presenciales con fundamentos conceptuales y
        análisis crítico de \textbf{casos reales} en logística,
        manufactura y salud.
      \end{conceptbox}
      \vspace{0.15cm}
      \begin{conceptbox}{\faFlask\enspace Laboratorios Interpretativos}
        \small
        Ejecutarán código Python \textbf{ya desarrollado} en
        Google Colab. El foco es \textbf{interpretar resultados},
        no programar desde cero.
      \end{conceptbox}
    \end{column}
    \begin{column}{0.48\textwidth}
      \begin{conceptbox}{\faSearch\enspace Análisis de Casos}
        \small
        Estudio de aplicaciones reales en marketing, operaciones
        y otros sectores industriales.
      \end{conceptbox}
      \vspace{0.15cm}
      \begin{conceptbox}{\faProjectDiagram\enspace Proyecto Final}
        \small
        Propuesta estratégica grupal basada en la interpretación
        de una solución de IA para un \textbf{problema industrial real}.
      \end{conceptbox}
    \end{column}
  \end{columns}
\end{frame}

% ─────────────────────────────────────────────────────────────
%  4. HERRAMIENTAS
% ─────────────────────────────────────────────────────────────
\begin{frame}{Ecosistema Tecnológico del Curso}
  \begin{columns}[T]
    \begin{column}{0.48\textwidth}
      \begin{block}{\faPython\enspace Lenguaje y Entorno}
        \small
        \begin{itemize}\setlength\itemsep{3pt}
          \item \textbf{Python} (uso interpretativo, no necesitan
                ser expertos).
          \item \textbf{Google Colab} / Jupyter Notebook como
                entorno en la nube.
        \end{itemize}
      \end{block}
    \end{column}
    \begin{column}{0.48\textwidth}
      \begin{block}{\faBoxes\enspace Librerías Principales}
        \small
        \begin{itemize}\setlength\itemsep{3pt}
          \item \tech{Pandas} y \tech{Numpy} --- Manejo de datos.
          \item \tech{Matplotlib} --- Visualización.
          \item \tech{Scikit-learn} --- Machine Learning.
          \item \tech{Keras} / \tech{TensorFlow} --- Deep Learning.
        \end{itemize}
      \end{block}
    \end{column}
  \end{columns}
  \vspace{0.3cm}
  \begin{alertblock}{\faLightbulb\enspace Tranquilos}
    \centering\small
    No necesitan instalar nada en su computador.
    Todo se ejecuta desde el \textbf{navegador} con Google Colab.
  \end{alertblock}
\end{frame}

% =============================================================
%  SECCIÓN 2: GOOGLE COLAB
% =============================================================

% ─────────────────────────────────────────────────────────────
%  5. ¿QUÉ ES GOOGLE COLAB?
% ─────────────────────────────────────────────────────────────
\begin{frame}{¿Qué es Google Colab?}
  \begin{columns}[T]
    \begin{column}{0.55\textwidth}
      \small
      Piensen en Google Colab como un \textbf{``cuaderno mágico''}
      que vive en su navegador:
      \begin{itemize}\setlength\itemsep{4pt}
        \item Permite ejecutar código \textbf{Python} sin instalar
              absolutamente nada.
        \item Es un entorno \textbf{gratuito} basado en Jupyter Notebook
              que vive en la nube de Google.
        \item Es su \textbf{laboratorio personal} para programar y
              analizar datos desde cualquier lugar.
      \end{itemize}
    \end{column}
    \begin{column}{0.42\textwidth}
      \begin{conceptbox}{\faCloud\enspace En Resumen}
        \small
        Google Colab = \\[3pt]
        \textbf{Navegador} + \textbf{Python} + \textbf{Nube}\\[3pt]
        Sin instalaciones. \\
        Sin configuraciones. \\
        Solo necesitan una cuenta de Google.
      \end{conceptbox}
    \end{column}
  \end{columns}
\end{frame}

% ─────────────────────────────────────────────────────────────
%  6. VENTAJAS DE COLAB
% ─────────────────────────────────────────────────────────────
\begin{frame}{Ventajas Clave de Google Colab}
  \begin{columns}[T]
    \begin{column}{0.48\textwidth}
      \begin{conceptbox}{\faLaptop\enspace Sin Instalación}
        \small
        Acceso directo desde el navegador. Nada de configuraciones
        complejas ni problemas de compatibilidad.
      \end{conceptbox}
      \vspace{0.15cm}
      \begin{conceptbox}{\faMicrochip\enspace GPU/TPU Gratuitas}
        \small
        Hardware potente sin costo para acelerar cálculos de
        Machine Learning cuando lo necesiten.
      \end{conceptbox}
    \end{column}
    \begin{column}{0.48\textwidth}
      \begin{conceptbox}{\faGlobe\enspace Accesible Globalmente}
        \small
        Trabajen desde cualquier dispositivo con conexión a internet.
        Solo necesitan un navegador.
      \end{conceptbox}
      \vspace{0.15cm}
      \begin{conceptbox}{\faGoogleDrive\enspace Integración con Drive}
        \small
        Guarden, organicen y compartan sus cuadernos a través de
        Google Drive de forma automática.
      \end{conceptbox}
    \end{column}
  \end{columns}
\end{frame}

% ─────────────────────────────────────────────────────────────
%  7. PRIMEROS PASOS
% ─────────────────────────────────────────────────────────────
\begin{frame}{Primeros Pasos en Google Colab}
  \vspace{0.2cm}
  \begin{center}
  \begin{tikzpicture}[font=\small,
    step/.style={draw=mainColor, fill=bgLight, rounded corners=4pt,
                 minimum width=3.2cm, minimum height=1.6cm,
                 align=center, text=codeGray}]
    \node[step] (s1) at (0,0)
      {\textbf{Paso 1}\\[2pt]\scriptsize Abrir\\
       \tech{colab.google.com}\\
       \scriptsize con cuenta Google};
    \node[step] (s2) at (5,0)
      {\textbf{Paso 2}\\[2pt]\scriptsize Menú \textit{Archivo}\\
       \scriptsize $\rightarrow$ \textit{Nuevo cuaderno}};
    \node[step] (s3) at (10,0)
      {\textbf{Paso 3}\\[2pt]\scriptsize Se guarda solo\\
       \scriptsize en \textit{Colab Notebooks}\\
       \scriptsize de Google Drive};
    \draw[-{Stealth}, thick, color=mainColor] (s1) -- (s2);
    \draw[-{Stealth}, thick, color=mainColor] (s2) -- (s3);
  \end{tikzpicture}
  \end{center}
  \vspace{0.3cm}
  \begin{alertblock}{\faLightbulb\enspace Tip}
    \centering\small
    Pueden acceder directamente desde \tech{colab.research.google.com}
    con su cuenta institucional o personal de Google.
  \end{alertblock}
\end{frame}

% ─────────────────────────────────────────────────────────────
%  8. ESTRUCTURA DE UN CUADERNO
% ─────────────────────────────────────────────────────────────
\begin{frame}[fragile]{Estructura de un Cuaderno: Celdas}
  \begin{columns}[T]
    \begin{column}{0.48\textwidth}
      \begin{conceptbox}{\faCode\enspace Celdas de Código}
        \small
        Aquí se escribe y ejecuta Python.\\[4pt]
        \textbf{¿Cómo ejecutar?}\\
        Botón \textit{Play} o \tech{Shift + Enter}\\[4pt]
        \textbf{Ejemplo:}
\begin{verbatim}
print("¡Hola, mundo!")
2 + 3
\end{verbatim}
      \end{conceptbox}
    \end{column}
    \begin{column}{0.48\textwidth}
      \begin{conceptbox}{\faFileAlt\enspace Celdas de Texto (Markdown)}
        \small
        Sirven para \textbf{documentar y explicar} el trabajo.\\[4pt]
        \textbf{Formato básico:}
\begin{verbatim}
# Título grande
## Subtítulo
**negrita** y *cursiva*
- Lista con viñetas
\end{verbatim}
      \end{conceptbox}
    \end{column}
  \end{columns}
  \vspace{0.2cm}
  \centering\small
  Un cuaderno de Colab alterna celdas de código y texto para crear
  un \textbf{documento interactivo} y reproducible.
\end{frame}

% ─────────────────────────────────────────────────────────────
%  9. PRIMER CÓDIGO
% ─────────────────────────────────────────────────────────────
\begin{frame}[fragile]{¡Hola, Mundo en la Nube!}
  \begin{block}{\faPython\enspace Tu primer código en Google Colab}
    \small Copien esto en una celda de código y presionen \tech{Shift + Enter}:
  \end{block}
  \vspace{0.2cm}
\begin{verbatim}
    # Esta es una celda de código en Google Colab
    print("¡Hola, mundo en Colab!")

    # Python como calculadora
    precio = 15000
    cantidad = 3
    total = precio * cantidad
    print(f"El total es: ${total}")
\end{verbatim}
  \vspace{0.2cm}
  \begin{alertblock}{\faCheckCircle\enspace Si ven el resultado...}
    \centering\small
    ¡Felicidades! Ya están ejecutando Python en la nube.
    No necesitaron instalar nada.
  \end{alertblock}
\end{frame}

% ─────────────────────────────────────────────────────────────
%  10. CONSEJOS PRÁCTICOS
% ─────────────────────────────────────────────────────────────
\begin{frame}{Consejos Prácticos para Google Colab}
  \begin{columns}[T]
    \begin{column}{0.32\textwidth}
      \begin{conceptbox}{\faTachometerAlt\enspace Runtime}
        \small
        Si el proyecto es pesado, activen la GPU en:\\[3pt]
        \textit{Entorno de ejecución} $\rightarrow$
        \textit{Cambiar tipo de entorno}.
      \end{conceptbox}
    \end{column}
    \begin{column}{0.32\textwidth}
      \begin{conceptbox}{\faSave\enspace Guardar Avances}
        \small
        Es automático, pero para versiones importantes usen:\\[3pt]
        \textit{Archivo} $\rightarrow$
        \textit{Guardar copia en Drive}.
      \end{conceptbox}
    \end{column}
    \begin{column}{0.32\textwidth}
      \begin{conceptbox}{\faCommentDots\enspace Comentar Código}
        \small
        Usen el símbolo \tech{\#} para que el código sea
        comprensible para ustedes y para otros.
      \end{conceptbox}
    \end{column}
  \end{columns}
  \vspace{0.3cm}
  \centering\small
  Recuerden: Colab se domina con la \textbf{práctica}.
  No tengan miedo de experimentar y equivocarse.
\end{frame}

% =============================================================
%  SECCIÓN 3: INTRODUCCIÓN A MACHINE LEARNING
% =============================================================

% ─────────────────────────────────────────────────────────────
%  11. ¿QUÉ ES MACHINE LEARNING?
% ─────────────────────────────────────────────────────────────
\begin{frame}{Introducción a Machine Learning}
  \begin{columns}[T]
    \begin{column}{0.55\textwidth}
      \small
      \textbf{Machine Learning} (Aprendizaje Automático) es una rama
      de la Inteligencia Artificial que permite a las máquinas
      \textbf{aprender a partir de datos} sin ser programadas
      explícitamente.\\[8pt]
      En lugar de instrucciones paso a paso, se usan
      \textbf{ejemplos} para que el sistema busque patrones
      y pueda \textbf{predecir o clasificar} información nueva.
    \end{column}
    \begin{column}{0.42\textwidth}
      \begin{conceptbox}{\faBrain\enspace Idea Central}
        \small
        \textbf{Programación clásica:}\\
        Datos + Reglas $\rightarrow$ Resultado\\[6pt]
        \textbf{Machine Learning:}\\
        Datos + Resultados $\rightarrow$ Reglas\\[6pt]
        La máquina \textit{descubre} las reglas
        a partir de los datos.
      \end{conceptbox}
    \end{column}
  \end{columns}
\end{frame}

% ─────────────────────────────────────────────────────────────
%  12. TIPOS DE APRENDIZAJE (CONTEXTO)
% ─────────────────────────────────────────────────────────────
\begin{frame}{Tipos de Aprendizaje Automático}
  \begin{columns}[T]
    \begin{column}{0.48\textwidth}
      \begin{conceptbox}{\faTag\enspace Supervisado}
        \small
        Le damos al modelo datos \textbf{con etiquetas}
        (la respuesta correcta).\\[4pt]
        \textbf{Ejemplos industriales:}
        \begin{itemize}\setlength\itemsep{2pt}
          \item Predecir demanda de productos.
          \item Clasificar piezas como defectuosas o no.
          \item Estimar precios de venta.
        \end{itemize}
      \end{conceptbox}
    \end{column}
    \begin{column}{0.48\textwidth}
      \begin{conceptbox}{\faObjectGroup\enspace No Supervisado}
        \small
        Los datos \textbf{no tienen etiquetas}. El modelo
        busca patrones por sí solo.\\[4pt]
        \textbf{Ejemplos industriales:}
        \begin{itemize}\setlength\itemsep{2pt}
          \item Segmentar clientes por comportamiento.
          \item Detectar anomalías en producción.
          \item Agrupar productos similares.
        \end{itemize}
      \end{conceptbox}
    \end{column}
  \end{columns}
\end{frame}

% ─────────────────────────────────────────────────────────────
%  13. ÁRBOLES DE DECISIÓN
% ─────────────────────────────────────────────────────────────
\begin{frame}[fragile]{Árboles de Decisión}
  \begin{columns}[T]
    \begin{column}{0.50\textwidth}
      \small
      Un árbol de decisión funciona como un \textbf{diagrama de flujo}:
      cada ``pregunta'' lleva a una rama distinta hasta llegar
      a una clasificación final.\\[8pt]
      \textbf{Ejemplo cotidiano:}
      \begin{enumerate}\setlength\itemsep{2pt}
        \item ¿Edad $>$ 18? \textbf{Sí}
        \item ¿Tiene licencia? \textbf{No}
        \item Resultado: \textcolor{accentColor}{\textbf{No puede alquilar auto}}
      \end{enumerate}
      \vspace{0.2cm}
      \textbf{En Python:}
\begin{verbatim}
from sklearn.tree import (
    DecisionTreeClassifier
)
model = DecisionTreeClassifier()
# model.fit(datos, etiquetas)
\end{verbatim}
    \end{column}
    \begin{column}{0.46\textwidth}
      \begin{center}
      \begin{tikzpicture}[
        font=\scriptsize,
        level distance=1.3cm,
        sibling distance=2.2cm,
        edge from parent/.style={draw, ->, thick, mainColor},
        every node/.style={draw=mainColor, fill=bgLight, rounded corners=3pt,
                          align=center, minimum width=1.8cm, inner sep=3pt}
      ]
        \node {¿Edad $>$ 18?}
          child { node {¿Licencia?}
            child { node[fill=mainColor!20] {\textbf{Sí puede}} edge from parent node[left, draw=none, fill=none] {Sí} }
            child { node[fill=accentColor!20] {\textbf{No puede}} edge from parent node[right, draw=none, fill=none] {No} }
            edge from parent node[left, draw=none, fill=none] {Sí}
          }
          child { node[fill=accentColor!20] {\textbf{No puede}} edge from parent node[right, draw=none, fill=none] {No} };
      \end{tikzpicture}
      \end{center}
    \end{column}
  \end{columns}
\end{frame}

% ─────────────────────────────────────────────────────────────
%  14. REGRESIÓN LINEAL
% ─────────────────────────────────────────────────────────────
\begin{frame}{Regresión Lineal}
  \begin{columns}[T]
    \begin{column}{0.55\textwidth}
      \small
      Método para predecir un \textbf{valor numérico continuo}
      encontrando la \textbf{mejor línea recta} que se ajuste
      a los datos.\\[8pt]
      \textbf{La ecuación:}
      \begin{equation*}
        y = mx + b
      \end{equation*}
      Donde:
      \begin{itemize}\setlength\itemsep{2pt}
        \item $y$ = valor predicho (ej: precio de venta).
        \item $x$ = variable de entrada (ej: metros cuadrados).
        \item $m$ = pendiente (cuánto cambia $y$ por cada unidad de $x$).
        \item $b$ = intersección (valor base).
      \end{itemize}
    \end{column}
    \begin{column}{0.42\textwidth}
      \begin{conceptbox}{\faChartLine\enspace Aplicaciones Industriales}
        \small
        \begin{itemize}\setlength\itemsep{3pt}
          \item Predecir \textbf{ventas futuras} a partir de datos históricos.
          \item Estimar \textbf{costos de producción}.
          \item Proyectar \textbf{demanda} de productos.
          \item Calcular \textbf{precio} de una propiedad según sus características.
        \end{itemize}
      \end{conceptbox}
    \end{column}
  \end{columns}
\end{frame}

% ─────────────────────────────────────────────────────────────
%  15. IA vs ML vs DL
% ─────────────────────────────────────────────────────────────
\begin{frame}{IA vs. Machine Learning vs. Deep Learning}
  \begin{center}
  \begin{tikzpicture}
    % Outer circle - IA
    \fill[mainColor!15] (0,0) circle (3.5cm);
    \draw[mainColor, thick] (0,0) circle (3.5cm);
    \node[font=\small\bfseries, text=mainColor] at (0, 3.0) {Inteligencia Artificial};
    \node[font=\scriptsize, text=codeGray, align=center] at (2.2, 2.2)
      {Sistemas que imitan\\la inteligencia humana};

    % Middle circle - ML
    \fill[accentColor!15] (0,-0.3) circle (2.3cm);
    \draw[accentColor, thick] (0,-0.3) circle (2.3cm);
    \node[font=\small\bfseries, text=accentColor] at (0, 1.6) {Machine Learning};
    \node[font=\scriptsize, text=codeGray, align=center] at (0, 0.9)
      {Aprende de datos\\sin programación explícita};

    % Inner circle - DL
    \fill[mainColor!30] (0,-0.8) circle (1.1cm);
    \draw[mainColor, thick] (0,-0.8) circle (1.1cm);
    \node[font=\small\bfseries, text=codeGray, align=center] at (0,-0.8)
      {Deep\\Learning};
  \end{tikzpicture}
  \end{center}
\end{frame}

% ─────────────────────────────────────────────────────────────
%  16. CIERRE
% ─────────────────────────────────────────────────────────────
\begin{frame}{¡El Viaje Apenas Comienza!}
  \begin{columns}[T]
    \begin{column}{0.55\textwidth}
      \small
      \textbf{Hoy aprendimos:}
      \begin{itemize}\setlength\itemsep{4pt}
        \item[\textcolor{mainColor}{\faCheck}] Qué es este curso y cómo se evalúa.
        \item[\textcolor{mainColor}{\faCheck}] Qué es Google Colab y cómo usarlo.
        \item[\textcolor{mainColor}{\faCheck}] Los conceptos básicos de Machine Learning.
        \item[\textcolor{mainColor}{\faCheck}] Qué son los árboles de decisión y la regresión lineal.
      \end{itemize}
      \vspace{0.3cm}
      \textbf{Próximas sesiones:}
      \begin{itemize}\setlength\itemsep{3pt}
        \item IA como ventaja competitiva.
        \item Librerías esenciales de Python.
        \item Manejo de datasets reales.
      \end{itemize}
    \end{column}
    \begin{column}{0.42\textwidth}
      \begin{alertblock}{\faLightbulb\enspace Para la próxima clase}
        \small
        \begin{enumerate}\setlength\itemsep{3pt}
          \item Crear cuenta en Google (si no tienen).
          \item Entrar a \tech{colab.research.google.com}.
          \item Crear un cuaderno nuevo y ejecutar:\\
                \tech{print("Listo")}
        \end{enumerate}
      \end{alertblock}
      \vspace{0.3cm}
      \centering
      \large\textbf{¡Prepárense para transformar\\datos en conocimiento!}
    \end{column}
  \end{columns}
\end{frame}

\end{document}
